\section{Interpret signed codes without corrupting their bytes}
The first left code is stored as \texttt{e0 2e}, representing 12000. The following
right code is \texttt{20 d1}. Combining those two unsigned bytes yields
\texttt{d120}, or 53536 as an unsigned number. The same sixteen-bit code represents
minus 12000 when interpreted as a signed two's-complement sample. Sign extension
belongs after assembly of the code, before multiplication by a signed coefficient.

A byte-oriented load path must not sign-extend \texttt{d1} as an independent component
and let its high bits contaminate another operation. Conversely, leaving the completed
code unsigned turns a negative sample into a large positive input. Those mistakes may
still produce valid WAV output with a plausible duration; an audio player is not a
signed-arithmetic checker.

A file's even sample offsets also do not prove alignment of the buffer address. If the
buffer begins at an odd address, its PCM samples remain odd-addressed after adding
an even offset such as 56. Use an input path that satisfies the callable kernel's
unaligned-byte contract. Test both the exact numerical samples and the bytes stored
for their results, including the distinct endpoints minus 32768 and 32767.
